В рамках данной работы была поставлена цель: предложить и разработать рабочий прототип альтернативного подхода к инференсу больших языковых моделей, основанных на архитектуре Mixture-of-Experts (MoE), удовлетворяющий требованиям динамической балансировки нагрузки и обеспечивающий более высокую отказоустойчивость по сравнению с оригинальным подходом, основанным на коллективах~\cite{shazeer2017moe,rajbhandari2022deepspeedmoe}.

Основная гипотеза заключалась в том, что замена коллективных All-to-All коммуникаций на уровне экспертного слоя в MoE моделях на клиент-серверное взаимодействие между основной частью модели и экспертными блоками позволяет удовлетворить поставленным требованиям. В качестве ориентира для коллективного подхода в работе рассматривается библиотека DeepEP~\cite{deepseekDeepEP}.

В процессе проверки гипотезы и проведения сопутствующих исследований был предложен и реализован прототип клиент-серверной архитектуры с асинхронной диспетчеризацией экспертов и механизмом передачи памяти между видеокартами. Система включает компоненты мониторинга состояния узлов, динамической балансировки нагрузки и обнаружения отказов в реальном времени.

Поскольку основная проблема выбранного подхода заключается в потенциальном снижении общей производительности, для ее решения были предложены и реализованы следующие оптимизации: уменьшение размера пакетов данных, которые передаются по сети, эффективное использование шины памяти за счет разбиения передачи данных на части, использование выравнивания экспертных активаций для более эффективного использования кэша GPU-операций, а также применение grouped GeMM для уменьшения количества точек синхронизации и повышения эффективности работы кэша GPU.

В результате проделанной работы спроектирован и разработан прототип экспертного слоя на основе клиент-серверной архитектуры для отказоустойчивого инференса MoE моделей, который обеспечивает возможность перераспределения нагрузки между экспертами и плавную деградацию производительности при отказах узлов, но демонстрирует производительность ниже, чем современные решения, основанные на коллективах. Система способна продолжать работу при выходе из строя отдельных экспертных блоков, в то время как решения, основанные на коллективных операциях, полностью останавливают инференс модели при возникновении отказа на любом из слоев.

В дальнейшей перспективе необходимы оптимизация асинхронного взаимодействия основной части модели и удалённых экспертных блоков, а также интеграция предложенного подхода с существующими системами инференса больших языковых моделей, такими как vLLM, SGLang и TensorRT-LLM~\cite{vllmProject,sglangProject,nvidiaTensorRTLLM}. Кроме того, представляет интерес разработка автоматических механизмов адаптации конфигурации системы под изменяющиеся паттерны нагрузки и создание библиотеки, реализующей предложенные техники для их широкого применения в индустрии.
